% Topic 1.3.05 — Numerical Differentiation, Integration, and ODEs.

\section{Численное диф., интегрирование и ОДУ}
\label{sec:1.3.05-numerical-diff-int-ode}

\begin{ticket}
Численное дифференцирование и интегрирование. Численные методы для решения систем дифференциальных уравнений.
\end{ticket}

\subsection*{Теория}

Рассматриваем три тесно связанные задачи приближения для
вещественнозначных функций одного вещественного аргумента: оценку
$f'(x)$ по значениям функции, оценку $\int_a^b f(x)\, dx$ по значениям
функции и приближённое решение задачи Коши для ОДУ. Всюду $h > 0$~---
шаг сетки. Подробное изложение см.\ в \cite[гл.~28]{CLRS2013} и
классических руководствах по вычислительной математике из главы~2.1.

\subsubsection*{Численное дифференцирование}

Заменяем предельное определение производной её приближением через
конечные разности.

\begin{definition}[Конечноразностные приближения]
\label{def:finite-diff}
Для достаточно гладкой~$f$ и шага $h > 0$:
\begin{align*}
  D_+ f(x) &= \frac{f(x + h) - f(x)}{h} && \text{разность вперёд,} \\
  D_- f(x) &= \frac{f(x) - f(x - h)}{h} && \text{разность назад,} \\
  D_0 f(x) &= \frac{f(x + h) - f(x - h)}{2h} && \text{центральная разность.}
\end{align*}
\end{definition}

\begin{theorem}[Погрешность конечноразностных приближений]
\label{thm:finite-diff-error}
Если $f \in C^2$ в окрестности~$x$, то
$D_+ f(x) - f'(x) = \tfrac{h}{2} f''(\xi_+)$ при некотором
$\xi_+ \in (x, x + h)$; в частности, $D_+ f(x) = f'(x) + O(h)$. То же
самое порядка~$h$ верно для~$D_-$. Если $f \in C^3$, то
$D_0 f(x) - f'(x) = \tfrac{h^2}{6} f'''(\xi_0)$ при некотором
$\xi_0 \in (x - h, x + h)$; в частности, $D_0 f(x) = f'(x) + O(h^2)$.
\end{theorem}
\proofof{finite-diff-error}

\subsubsection*{Численное интегрирование (квадратуры)}

\emph{Квадратурная формула} приближает $\int_a^b f(x)\, dx$ конечной
линейной комбинацией значений~$f$ в точках. Простейшие формулы
получаются интерполяцией~$f$ многочленом по равноотстоящим узлам и
точным интегрированием интерполянта. Разбивая $[a, b]$ на $n$
подотрезков длины $h = (b - a)/n$ с узлами $x_k = a + kh$,
$k = 0, \dots, n$, получают \emph{составные} версии формул.

\begin{definition}[Составные квадратурные формулы]
\label{def:composite-quad}
При $f_k = f(x_k)$:
\begin{itemize}
  \item составная \emph{формула средних}:
        $M_n(f) = h \sum_{k=0}^{n-1} f\!\bigl(x_k + h/2\bigr)$;
  \item составная \emph{формула трапеций}:
        $T_n(f) = h \bigl(\tfrac{f_0}{2} + f_1 + f_2 + \dots + f_{n-1} + \tfrac{f_n}{2}\bigr)$;
  \item составная \emph{формула Симпсона}: при чётном~$n$
        $S_n(f) = \tfrac{h}{3}\bigl(f_0 + 4 f_1 + 2 f_2 + 4 f_3 + \dots + 2 f_{n-2} + 4 f_{n-1} + f_n\bigr)$.
\end{itemize}
\end{definition}

\begin{theorem}[Оценки погрешности составных формул]
\label{thm:quadrature-error}
Пусть $I(f) = \int_a^b f(x)\, dx$ и $h = (b - a)/n$.
\begin{enumerate}[label=(\alph*)]
  \item Если $f \in C^2[a, b]$, то
        $|T_n(f) - I(f)| \leq \tfrac{(b - a) h^2}{12} \max_{[a, b]} |f''|$;
        в частности, $T_n(f) = I(f) + O(h^2)$.
  \item Если $f \in C^4[a, b]$ и $n$ чётно, то
        $|S_n(f) - I(f)| \leq \tfrac{(b - a) h^4}{180} \max_{[a, b]} |f^{(4)}|$;
        в частности, $S_n(f) = I(f) + O(h^4)$.
  \item Формула Симпсона \emph{точна} (даёт нулевую погрешность) на
        многочленах степени не выше~$3$.
\end{enumerate}
\end{theorem}
\proofof{quadrature-error}

\subsubsection*{Задача Коши для ОДУ}

Рассмотрим задачу с начальным условием (задачу Коши)
\[
  y'(x) = f(x, y(x)),\qquad y(x_0) = y_0,
\]
на отрезке $[x_0, x_0 + L]$, при условии, что~$f$ липшицева по~$y$
равномерно по~$x$ (и значит, решение существует и единственно по
теореме Пикара--Линделёфа). Численный метод выдаёт приближённые
значения $y_n \approx y(x_n)$ на сетке $x_n = x_0 + n h$,
$n = 0, 1, \dots, N$, при $N h = L$.

\begin{definition}[Явный метод Эйлера]
\label{def:euler}
$y_{n+1} = y_n + h\, f(x_n, y_n)$, $n = 0, 1, \dots, N - 1$.
\end{definition}

\begin{definition}[Неявный метод Эйлера]
\label{def:implicit-euler}
$y_{n+1} = y_n + h\, f(x_{n+1}, y_{n+1})$~--- на каждом шаге
неизвестный $y_{n+1}$ находится из неявного уравнения.
\end{definition}

\begin{definition}[Классический метод Рунге--Кутта (RK4)]
\label{def:rk4}
\[
  k_1 = f(x_n, y_n),\quad
  k_2 = f\!\bigl(x_n + \tfrac{h}{2},\, y_n + \tfrac{h}{2} k_1\bigr),\quad
  k_3 = f\!\bigl(x_n + \tfrac{h}{2},\, y_n + \tfrac{h}{2} k_2\bigr),\quad
  k_4 = f(x_n + h,\, y_n + h k_3),
\]
\[
  y_{n+1} = y_n + \tfrac{h}{6}\bigl(k_1 + 2 k_2 + 2 k_3 + k_4\bigr).
\]
\end{definition}

\begin{theorem}[Сходимость явного метода Эйлера]
\label{thm:euler-convergence}
Пусть $f$ $L$-липшицева по~$y$ (равномерно по~$x$) в рассматриваемой
полосе, точное решение $y \in C^2[x_0, x_0 + L]$ и
$\|y''\|_{\infty} \leq M$. Тогда явный метод Эйлера для всякого~$n$
с $x_n \leq x_0 + L$ удовлетворяет
\[
  |y_n - y(x_n)| \;\leq\; \frac{M h}{2 L} \bigl(e^{L (x_n - x_0)} - 1\bigr) \;=\; O(h).
\]
В частности, метод Эйлера имеет порядок сходимости~$1$.
\end{theorem}
\proofof{euler-convergence}

\begin{remark}[Словарь: порядок, согласованность, устойчивость]
Одношаговый метод имеет \emph{локальную ошибку усечения} $\tau_n$
порядка $h^{p+1}$, если $\tau_n = O(h^{p+1})$, и \emph{глобальную
ошибку} порядка $O(h^p)$ при наличии устойчивости.
\emph{Согласованность} ($\tau_n \to 0$ при $h \to 0$) вместе с
\emph{нуль-устойчивостью} (малые возмущения остаются ограниченными)
дают \emph{сходимость}~--- это теорема эквивалентности Лакса в данной
постановке. Метод Эйлера имеет порядок~$1$, RK4~--- порядок~$4$:
существенно меньше вызовов~$f$ при той же точности на гладких задачах.
\end{remark}

\subsection*{Доказательства}

\begin{proof}[\Cref{thm:finite-diff-error}]
\proofanchor{finite-diff-error}
\emph{Разность вперёд.} Разложение Тейлора с остатком в форме Лагранжа
(\cref{thm:taylor-lagrange}, применённое к~$f$ в окрестности~$x$):
\[
  f(x + h) = f(x) + h f'(x) + \tfrac{h^2}{2} f''(\xi_+),
  \quad \xi_+ \in (x, x + h).
\]
Решая относительно $D_+ f(x) = (f(x + h) - f(x))/h$, получаем
$D_+ f(x) = f'(x) + \tfrac{h}{2} f''(\xi_+)$, откуда
$D_+ f(x) - f'(x) = \tfrac{h}{2} f''(\xi_+)$~--- порядка~$h$.

\emph{Центральная разность.} Два разложения Тейлора с членом третьего
порядка:
\begin{align*}
  f(x + h) &= f(x) + h f'(x) + \tfrac{h^2}{2} f''(x) + \tfrac{h^3}{6} f'''(\xi_1), \\
  f(x - h) &= f(x) - h f'(x) + \tfrac{h^2}{2} f''(x) - \tfrac{h^3}{6} f'''(\xi_2),
\end{align*}
где $\xi_1 \in (x, x + h)$, $\xi_2 \in (x - h, x)$. Вычитая и деля на
$2h$, получаем
$D_0 f(x) = f'(x) + \tfrac{h^2}{12}(f'''(\xi_1) + f'''(\xi_2))$. По
свойству промежуточного значения для $f'''$ на $(x - h, x + h)$
среднее $(f'''(\xi_1) + f'''(\xi_2))/2$ равно $f'''(\xi_0)$ при
некотором $\xi_0$ из этого интервала, что даёт
$D_0 f(x) - f'(x) = \tfrac{h^2}{6} f'''(\xi_0)$.
\end{proof}

\begin{proof}[\Cref{thm:quadrature-error}]
\proofanchor{quadrature-error}
\emph{(а) Трапеции.} На отдельном подотрезке $[x_k, x_{k+1}]$ длины~$h$
формула трапеций $\tfrac{h}{2}(f_k + f_{k+1})$ совпадает с
$\int_{x_k}^{x_{k+1}} L_1(x)\, dx$ для линейного интерполянта~$L_1$,
проходящего через $(x_k, f_k), (x_{k+1}, f_{k+1})$. Формула остатка
интерполяции (Лагранж для одноточечной интерполяции типа Эрмита)
даёт $f(x) - L_1(x) = \tfrac{(x - x_k)(x - x_{k+1})}{2}\, f''(\eta(x))$.
Интегрируя по $[x_k, x_{k+1}]$: интеграл от
$(x - x_k)(x - x_{k+1})$ равен $-h^3/6$, откуда
\[
  \Bigl| \int_{x_k}^{x_{k+1}} f(x)\, dx - \tfrac{h}{2}(f_k + f_{k+1}) \Bigr|
  \leq \frac{h^3}{12} \max_{[x_k, x_{k+1}]} |f''|.
\]
Суммируя по $n$ подотрезкам и используя $n h = b - a$, получаем
суммарную оценку $\tfrac{(b - a) h^2}{12} \max |f''|$.

\emph{(б) Симпсон.} Одиночная формула Симпсона на $[x_k, x_{k+2}]$~---
интеграл квадратичного интерполянта по точкам
$(x_k, f_k), (x_{k+1}, f_{k+1}), (x_{k+2}, f_{k+2})$. Аналогичное
разложение Тейлора в окрестности~$x_{k+1}$ (или представление через
ядро Пеано) даёт погрешность на одном блоке
$\tfrac{h^5}{90} f^{(4)}(\eta_k)$ при некотором
$\eta_k \in (x_k, x_{k+2})$; суммируя по $n/2$ блокам, получаем
$\tfrac{(b - a) h^4}{180} \max |f^{(4)}|$.

\emph{(в) Точность на кубических многочленах.} И квадратичные, и
кубические интерполянты Симпсон интегрирует точно, поскольку главный
член погрешности содержит $f^{(4)}$, обращающееся в нуль для
многочленов степени $\leq 3$. Это специфическая алгебраическая
особенность~--- формула Симпсона имеет \emph{степень точности}~$3$,
несмотря на то что построена на квадратичной интерполяции.
\end{proof}

\begin{proof}[\Cref{thm:euler-convergence}]
\proofanchor{euler-convergence}
\emph{Локальная ошибка усечения.} По формуле Тейлора для точного
решения:
$y(x_{n+1}) = y(x_n) + h y'(x_n) + \tfrac{h^2}{2} y''(\xi_n)
            = y(x_n) + h f(x_n, y(x_n)) + \tfrac{h^2}{2} y''(\xi_n)$.
Вычитая шаг Эйлера $y_{n+1} = y_n + h f(x_n, y_n)$:
\[
  e_{n+1} := y_{n+1} - y(x_{n+1}) = e_n + h \bigl(f(x_n, y_n) - f(x_n, y(x_n))\bigr) - \tfrac{h^2}{2} y''(\xi_n).
\]
По условию Липшица
$|f(x_n, y_n) - f(x_n, y(x_n))| \leq L |e_n|$, откуда
$|e_{n+1}| \leq (1 + h L) |e_n| + \tfrac{M h^2}{2}$.

\emph{Глобальная ошибка.} Итерируем с $e_0 = 0$:
$|e_n| \leq \tfrac{M h^2}{2} \sum_{k=0}^{n-1} (1 + h L)^k
        = \tfrac{M h^2}{2} \cdot \tfrac{(1 + h L)^n - 1}{h L}
        = \tfrac{M h}{2 L} \bigl((1 + h L)^n - 1\bigr)$.
Поскольку $(1 + h L)^n \leq e^{n h L} = e^{L (x_n - x_0)}$, получаем
$\tfrac{M h}{2 L} (e^{L (x_n - x_0)} - 1)$. Множитель при~$h$
ограничен при $x_n \leq x_0 + L$, так что $|e_n| = O(h)$.
\end{proof}

\subsection*{Примеры}

\begin{example}[Разность вперёд против центральной для $\sin'(0)$]
\label{ex:diff-sin}
Точное значение: $\sin'(0) = \cos(0) = 1$. При $h = 0{,}1$:
\begin{itemize}
  \item вперёд: $D_+ \sin(0) = (\sin(0{,}1) - 0)/0{,}1 \approx 0{,}99833$;
        погрешность $\approx 1{,}7 \times 10^{-3}$, что согласуется с
        оценкой $h/2$;
  \item центральная: $D_0 \sin(0) = (\sin(0{,}1) - \sin(-0{,}1))/0{,}2$;
        ввиду нечётности синуса числитель равен $2 \sin(0{,}1)$ и
        $D_0 \sin(0) = \sin(0{,}1)/0{,}1 \approx 0{,}99833$~--- так же,
        как и для $D_+$. Это совпадение возникает потому, что в
        точке~$0$ центральная формула для нечётной функции вырождается.
        Возьмём вместо этого $\sin'(\pi/3) = 0{,}5$: тогда
        $D_+ \sin(\pi/3) \approx 0{,}45524$ (погрешность
        $\approx 4{,}5 \times 10^{-2}$), а
        $D_0 \sin(\pi/3) \approx 0{,}49917$ (погрешность
        $\approx 8{,}3 \times 10^{-4}$). Центральная формула выигрывает
        на два порядка по~$h$~--- как и предсказано.
\end{itemize}
\end{example}

\begin{example}[Трапеции против Симпсона для $\int_0^1 e^x\, dx$]
\label{ex:quadrature-exp}
Точное значение $e - 1 \approx 1{,}71828$.

При $n = 4$, $h = 0{,}25$: формула трапеций даёт
$T_4 \approx 1{,}72722$ (погрешность $\approx 8{,}9 \times 10^{-3}$);
формула Симпсона~--- $S_4 \approx 1{,}71832$ (погрешность
$\approx 3{,}6 \times 10^{-5}$). При уменьшении $h$ вдвое до $h = 0{,}125$
($n = 8$) погрешность трапеций падает до $\approx 2{,}2 \times 10^{-3}$
(в~$4$ раза, согласно $h^2$); погрешность Симпсона~--- до
$\approx 2{,}2 \times 10^{-6}$ (в~$16$ раз, согласно $h^4$). Каждая
формула подтверждает заявленную в \cref{thm:quadrature-error}
скорость сходимости.
\end{example}

\begin{example}[Метод Эйлера для модельной задачи $y' = -y$]
\label{ex:euler-model}
$y' = -y$, $y(0) = 1$, точное решение $y(x) = e^{-x}$. Шаг Эйлера~---
$y_{n+1} = y_n - h y_n = (1 - h) y_n$, откуда $y_n = (1 - h)^n$. При
$x = 1$, $h = 0{,}1$, $n = 10$: $y_{10} = 0{,}9^{10} \approx 0{,}34868$
против $e^{-1} \approx 0{,}36788$; погрешность
$\approx 1{,}9 \times 10^{-2}$. Уменьшая $h$ вдвое до $0{,}05$:
$y_{20} = 0{,}95^{20} \approx 0{,}35849$, погрешность
$\approx 9{,}4 \times 10^{-3}$~--- примерно в~$2$ раза меньше, что
согласуется с порядком сходимости~$1$. При $h \geq 2$ имеем
$|1 - h| \geq 1$, и $y_n$ не убывает: метод Эйлера на этой задаче
\emph{условно устойчив} (с ограничением $h < 2/|\lambda|$ для
уравнения $y' = \lambda y$).
\end{example}

\begin{problem}[Степень точности формулы Симпсона]
Проверьте прямым вычислением, что на одном подотрезке $[a, b]$ длины
$h = b - a$ простая формула Симпсона
$\tfrac{h}{6}\bigl(f(a) + 4 f((a + b)/2) + f(b)\bigr)$ точно
интегрирует одночлены $1, x, x^2, x^3$, а на $x^4$ ошибается на
$\tfrac{h^5}{120}$ (вычислите оба интеграла и возьмите разность).
Используйте это, чтобы вывести главный член погрешности порядка
$h^4$ из \cref{thm:quadrature-error}~(б).
\end{problem}
